% ============================================================
% NIDA Akademik Rapor - Türkçe Özet
% ============================================================

\chapter*{ÖZET}
\addcontentsline{toc}{chapter}{ÖZET}

Bu çalışma, psikolojik danışmanlık ve rehberlik (PDR) alanında çocukların psikolojik durumlarının değerlendirilmesinde kullanılan projektif çizim testlerinin, yapay zekâ teknolojileri ile desteklenmesi amacıyla geliştirilmiş bir dijital uygulama projesini kapsamaktadır. \textbf{NIDA (Neural Imagery Decision Assessment)}, 5-12 yaş arası çocukların House-Tree-Person (HTP) çizimlerini analiz ederek okul psikolojik danışmanlarına ve rehber öğretmenlere akademik literatüre dayalı yorumlar sunan web tabanlı bir karar destek sistemidir.

Çalışmanın temel amacı, okul ortamlarında psikolojik danışmanların çocuk çizimlerini değerlendirme süreçlerini desteklemek, erken dönem psikolojik sorunların tespitinde yardımcı bir araç sunmak ve PDR uygulamalarının teknoloji ile entegrasyonuna katkı sağlamaktır. Sistem, John Buck (1948) tarafından geliştirilen HTP testinin yorumlama kurallarını, güncel akademik literatür ile harmanlayarak yapılandırılmış bir bilgi tabanı (Knowledge Base) oluşturmuş ve bu bilgi tabanını çok modlu büyük dil modelleri (Multimodal Large Language Models - MLLM) ile entegre etmiştir.

NIDA, tanı koyan bir araç olmayıp, profesyonel değerlendirmeyi destekleyen bir karar destek sistemi olarak konumlandırılmıştır. Sistem, her yorumu akademik kaynaklara dayandırmakta, güven seviyelerini belirtmekte ve çelişkili akademik görüşleri şeffaf biçimde sunmaktadır. Bu yaklaşım, PDR alanının etik ilkeleri ve bilimsel temelleri ile uyumlu bir teknoloji entegrasyonu sağlamaktadır.

Çalışma, eğitim teknolojilerinin psikolojik danışmanlık hizmetlerine entegrasyonu, yapay zekânın PDR alanındaki potansiyel kullanımları ve çocuk değerlendirmesinde dijital araçların rolü konularında alanyazına katkı sunmaktadır.

\vspace{1cm}

\textbf{Anahtar Kelimeler:} Psikolojik danışmanlık ve rehberlik, House-Tree-Person testi, projektif çizim testleri, yapay zekâ, çok modlu dil modelleri, eğitim teknolojisi, çocuk değerlendirmesi

\newpage
